\documentclass{article}
\usepackage[utf8]{inputenc}
\usepackage{amsmath}
\usepackage{geometry}
\geometry{margin=2cm}
\begin{document}

\section*{Análise da Questão 168 — ENEM 2024 — Matemática}

Esta questão apresenta um experimento que mede o comportamento oscilatório de amortecedores em duas superfícies distintas, retratado por funções com comportamento senoidal. O problema destaca duas condições para aprovação: as amplitudes nas duas superfícies (asfalto e estrada de chão) devem ser constantes, e a amplitude no asfalto deve ser menor do que na estrada de chão.

\subsection*{Solução Técnica}
A alternativa correta é a \textbf{A}, dado que apresenta um gráfico com amplitudes constantes em ambos os trechos, e que a amplitude durante o período em asfalto (0 a 9 minutos) é menor que a amplitude na estrada de chão (9 a 20 minutos).

\subsection*{Formulação e passos para análise}
\begin{enumerate}
  \item Identificar os intervalos temporais do gráfico: asfalto $(0,9)$ e estrada de chão $(9,20)$ minutos.
  \item Verificar a constância da amplitude da senoide em cada intervalo.
  \item Comparar as amplitudes dos dois intervalos para validar a condição de amplitude menor no asfalto.
\end{enumerate}

\subsection*{Fórmula chave}
\[
\text{Amplitude} = \max(|y|)\quad \text{(distância máxima vertical da curva em relação ao eixo zero)}
\]

\subsection*{Evidências visuais}
\begin{itemize}
  \item Gráfico da alternativa A mostra claramente duas seções senoidais distintas, cada uma com amplitude constante e a do asfalto menor.
  \item Comparação entre amplitudes nos períodos 0-9 e 9-20 minutos evidenciam a diferença que justifica a aprovação.
\end{itemize}

\subsection*{Explicação didática resumida}
Deve-se observar que o gráfico ideal demonstra oscilações estáveis (amplitude constante) em cada intervalo temporal e que a intensidade da oscilação (amplitude) é menor no asfalto. Isto significa que o amortecedor está funcionando adequadamente, moderando as oscilações na parte asfaltada e permitindo maior amplitude na estrada de chão.

\subsection*{Erros comuns e dicas didáticas}
\begin{itemize}
  \item Confundir amplitude com frequência das oscilações.
  \item Escolher gráficos com amplitudes variáveis por não observarem a constância requerida.
  \item Ignorar a diferença requerida entre as amplitudes nos dois trechos.
  \item Explique para os alunos que a amplitude representa o "alcance" da oscilação e que valores constantes indicam estabilidade na resposta do amortecedor.
\end{itemize}

\subsection*{Distratores e suas possíveis interpretações}
\begin{itemize}
  \item \textbf{B}: Apresenta amplitudes constantes, mas com a amplitude no asfalto maior que na estrada de chão, invertendo a condição exigida.
  \item \textbf{C}: Apresenta amplitude constante porém igual nos dois trechos, não atendendo à condição de amplitude menor no asfalto.
  \item \textbf{D}: Mostra amplitudes constantes, mas a amplitude no asfalto não é menor, confundindo com frequência erradamente.
  \item \textbf{E}: Apresenta amplitude constante no asfalto, porém no trecho da estrada de chão a amplitude não é constante, violando a condição para aprovação.
\end{itemize}

\subsection*{Perfil da questão}
A questão avalia a interpretação visual de gráficos matemáticos, especificamente funções senoides, e a compreensão dos conceitos de amplitude e constância da amplitude. Exige habilidade média e interpretação visual precisa, além de raciocínio lógico para comparar os intervalos do gráfico e extrair conclusões coerentes.

% Imagens das alternativas (representações textuais indicam os gráficos conforme fornecido no enunciado e análise)
\subsection*{Representação dos gráficos}

% A representar de forma textual resumida para fins didáticos:
\begin{itemize}
  \item \textbf{Alternativa A}: Amplitudes constantes, com amplitude menor no intervalo 0-9 minutos.
  \item \textbf{Alternativa B}: Amplitudes constantes, mas amplitude maior no asfalto.
  \item \textbf{Alternativa C}: Amplitudes constantes e iguais nos dois intervalos.
  \item \textbf{Alternativa D}: Amplitudes constantes, mas amplitude no asfalto maior ou igual que na estrada.
  \item \textbf{Alternativa E}: Amplitude constante no asfalto, porém variável na estrada de chão.
\end{itemize}

Em resumo, a alternativa A é a única que satisfaz as condições propostas pelo enunciado para a aprovação do amortecedor no experimento.

\vspace{1em}
\begin{center}
\textit{Análise completa e detalhada para entendimento da questão e seu contexto.}
\end{center}


\end{document}